\documentclass[10pt,a4paper]{amsart}
\usepackage[margin=19mm]{geometry}
\usepackage{amsmath,amssymb,mathtools}
\usepackage[scr=rsfs,cal=euler]{mathalfa}
\usepackage{xfrac}
\usepackage[unicode,hidelinks,bookmarks=false]{hyperref}
\newcommand{\ie}{i.e.\ }
\newcommand{\cf}{cf.\ }
\newcommand{\resp}{respectively\ }
\newcommand{\Subsection}{\subsection}
\providecommand{\nobreakdash}{\nobreak\hbox{-}}
\setlength{\emergencystretch}{2em}
\multlinegap0pt
\usepackage{bm}
\usepackage[normalem]{ulem}
\usepackage{siunitx}
\usepackage{mathtools}
\usepackage{amssymb}
\usepackage{float}
\usepackage{xcolor}
\usepackage{algorithm}\usepackage{algpseudocode}
\newtheorem{theorem}{Theorem}
\newcommand{\R}{\mathbb{R}}                                 % Real numbers
\newcommand{\pmat}[1]{\begin{pmatrix}#1\end{pmatrix}}
\title{Model order reduction via Lie groups}
\author{Yannik P. Wotte, Patrick Buchfink, Silke Glas, Federico Califano, Stefano Stramigioli}
\date{Source: arXiv:2511.03520v2 (CC BY 4.0)}
\begin{document}
\maketitle

\noindent\emph{Source and licence.} Model order reduction via Lie groups. Version arXiv:2511.03520v2 (CC BY 4.0), distributed by arXiv under the Creative Commons Attribution 4.0 International licence, http://creativecommons.org/licenses/by/4.0/, as recorded in the arXiv licence metadata for this exact version and shown on the abstract page https://arxiv.org/abs/2511.03520v2. Full licence notice: \texttt{licenses/arxiv-2511.03520-CC-BY-4.0.txt}.\par\smallskip

\noindent\textit{Editorial scope.} Counted unit: the article's theorem theorem (source0.tex lines 2577--2588). The article's complete original proof of it is delivered with it (source0.tex lines 2590--2735, one \texttt{proof} environment of 146 lines). No mathematics is written, repaired or re-derived: the statement and the proof are the article's own lines, in the article's own notation, and no retained line is substituted or re-flowed.  Retained uncounted as the source-local context the proof needs: lines 1419--1421; lines 1438--1440.  Nothing was omitted from any retained span, so the package carries no omission record.  No retained line contains a citation, so the wrapper carries no bibliography and no bibliography entry.  No retained line contains a figure environment, an image-inclusion command or drawing code, so no image is delivered and the package declares no asset dependency.  Editorial formatting only, in addition: an AMS \texttt{amsart} preamble that restates the article's own declarations and macros used by the retained lines, namely the environments \texttt{theorem} on separate counters, together with the standard packages those lines need.  The retained environments print in this document as Theorem 1 in order of appearance, whereas the article prints the same items with the numbering of its own full document.  The complete native source of the article as distributed in the arXiv e-print for this exact version is bundled unchanged as \texttt{sources/188514/source0.tex}.
\begin{equation}\label{eq:rigid-group-dynamics}
    \dot{H}(t) = H(t)\widetilde{T}(t)\,,\; H(0) = H_0\,.
\end{equation}

\begin{equation}\label{eq:snapshots-rigid-pointcloud}
    S := \bigg[\pmat{P_{0,0} \\ t_0}, \cdots, \pmat{P_{N_j,N_k}\\ t_{N_k}}\bigg] \,.
\end{equation}  

\begin{theorem}
    Assume the following:
    \begin{enumerate}
        \item The number of columns $N_jN_k$ and the number of rows $3N_i+1$ of $S$ are sufficiently large: $N_jN_k \geq 9(N_j+1)+3$ and $3N_i+1 \geq 9(N_j+1)+3$
        %\item The noise is $\eta^i_{j,k} = 0$
        \item Each initial point cloud $P_{j,0} \in \R^{N_j}$ spans $\R^3$: $$\dim\textnormal{span}\{p^i_{j,0}\;|\;i \in \{0,\cdots,N_i\}\} = 3$$
        \item The matrices $R(t_k)$ span $\R^{3\times 3}$: $$\dim\textnormal{span}\{R(t_k)\;|\;k \in \{0,\cdots,N_k\}\} = 9$$
        \item The translations $b(t_k)$ span $\R^3$: $$\dim\textnormal{span}\{b(t_k)\;|\;k \in \{0,\cdots,N_k\}\} = 3$$
        \item The subspaces $$W_j := \textnormal{span}\{P_{j,k}\;|\; k \in \{0,\cdots,N_k\} \}$$ of $\R^{3N_i}$ associated with different point cloud trajectories $P_{j,k}$ are independent. 
    \end{enumerate}  
    Then the number of non-zero singular values of the snapshot set $S$ in~\eqref{eq:snapshots-rigid-pointcloud} is $9(N_j+1)+3$.
\end{theorem}

\begin{proof}
    The number of non-zero singular values of a matrix is bounded by the number of columns $N_jN_k$, the number of rows $3N_i+1$, and the column rank of $S$. Given Assumption 1. that $N_jN_k$ and $3N_i+1$ are sufficiently large, we show that the column rank of $S$ is $9(N_j+1)+3$.
    
    Each column of $S$ corresponding to trajectory $j$ at time $t_k$ has the form
    \begin{equation}
    P_{j,k}     =
    \begin{pmatrix}
    R(t_k)p^1_{j,0} \\ 
    \vdots \\
    R(t_k)p^{N_i}_{j,0}
    \end{pmatrix}
    +
    \begin{pmatrix}
    b(t_k) \\
    \vdots \\
    b(t_k)
    \end{pmatrix}.
    \end{equation}
    We analyze rotational and translational contributions separately, first focusing on the rotational term.

    Write as $e_n \in \R^3$ the $n$-th unit vector, and define vectors $\bar P^{(j)}_{nm} \in \mathbb{R}^{3N_i}$ by
    \begin{equation}
        \bar P^{(j)}_{nm}
        =
        \begin{pmatrix}
        (e_n^\top p^1_{j,0}) e_m \\
        \vdots \\
        (e_n^\top p^{N_i}_{j,0}) e_m
        \end{pmatrix}\,.
    \end{equation}
    Then
    \begin{align}
        P^{\mathrm{rot}}_{j,k} = 
            \pmat{
            R(t_k)p^1_{j,0} \\ 
            \vdots \\
            R(t_k)p^{N_i}_{j,0} }
        &= \sum_{n,m} R_{nm}(t_k)
            \pmat{ 
            (e_n^\top p^1_{j,0}) e_m \\
            \vdots \\
            (e_n^\top p^{N_i}_{j,0}) e_m } \\ 
        &= \sum_{n,m} R_{nm}(t_k) \bar{P}^{(j)}_{nm} \,.
    \end{align}
    Hence the rotational part lies in the span of at most $9$ vectors. We now show that these are linearly independent. Suppose that
    \begin{equation}
    \sum_{n,m=1}^{3} a_{nm}\bar P^{(j)}_{nm}=0.
    \end{equation}
    Then for each particle $p^i_{j,0}$
    \begin{equation}
    \sum_{n,m} a_{nm}(e_n^\top p^i_{j,0}) e_m = 0\,.
    \end{equation}
    Define
    \begin{equation}
    A := \sum_{n,m} a_{nm} e_m e_n^\top \in \mathbb{R}^{3\times3}.
    \end{equation}
    The previous equation is equivalent to
    \begin{equation}
    A p^i_{j,0} = 0 \quad \text{for all } i.
    \end{equation}
    By Assumption 2., the initial cloud spans $\mathbb{R}^3$, so this would imply $A=0$, and therefore $a_{nm}=0$ for all $n,m$. Thus, the nine vectors $\bar P^{(j)}_{nm}$ are linearly independent.

    By Assumption 3., the rotational component for trajectory $j$ spans a $9$-dimensional subspace of $\mathbb{R}^{3N_i}$.

    By Assumption 5., it follows that $9$ linearly independent components are required for the rotational component of each of the $N_j + 1$ trajectories, and we arrive at $9(N_j+1)$ linearly independent components to express the rotational component of $S$. 
    
    By Assumption 4., the translations can be expressed by three linearly independent vectors $\bar{P}_{b,1},\cdots, \bar{P}_{b,3}$ for all trajectories, bringing the maximum number of orthogonal components to $9(N_j+1)+3$. 
    % \medskip
    % \noindent
    % \textbf{Step 3: Multiple trajectories.}

    % For distinct trajectories $j$, the associated subspaces
    % \begin{equation}
    % \mathcal V_j := \operatorname{span}\{\bar P^{(j)}_{nm}\}_{n,m=1}^{3}
    % \end{equation}
    % are generically independent provided the initial point clouds are in generic relative position. Hence, generically,
    % \begin{equation}
    % \dim \operatorname{span}\bigcup_{j=0}^{N_j} \mathcal V_j
    % =
    % 9(N_j+1).
    % \end{equation}

    % \medskip
    % \noindent
    % \textbf{Step 4: Translational contribution.}

    % The translation part has the form
    % \begin{equation}
    % P^{\mathrm{trans}}_{j,k}
    % =
    % \begin{pmatrix}
    % b(t_k) \\
    % \vdots \\
    % b(t_k)
    % \end{pmatrix}.
    % \end{equation}
    % Define
    % \begin{equation}
    % \bar P_{b,\ell}
    % =
    % \begin{pmatrix}
    % e_\ell \\
    % \vdots \\
    % e_\ell
    % \end{pmatrix},
    % \qquad \ell=1,2,3.
    % \end{equation}
    % If $\dim \operatorname{span}\{b(t_k)\}_{k=0}^{N_k}=3$, then the translation vectors span a $3$-dimensional subspace.

    % Moreover, these vectors are generically independent of the rotational subspaces, since rotational components scale with particle-dependent coefficients $(e_n^\top p^i_{j,0})$, whereas translations are constant across particles.

    % Thus translations contribute $3$ additional independent directions.

    % \medskip
    % \noindent
    % \textbf{Step 5: Rank computation.}

    % Under the stated dimension assumptions,
    % \begin{equation}
    % \dim \operatorname{span}\{P_{j,k}\}
    % =
    % 9(N_j+1)+3.
    % \end{equation}
    % Finally, the rank of $S$ is bounded by the number of columns $N_jN_k$ and the number of rows $3N_i+1$. Therefore,
    % \begin{equation}
    % \operatorname{rank}(S)
    % =
    % \min\big(9(N_j+1)+3,\; N_jN_k,\; 3N_i+1\big),
    % \end{equation}
    % which holds generically.

    % By definition, the snapshot set $S$ consists of combined rotations and translations of static point clouds $P \in \mathbb{R}^{3N_i}$. Identify $H \in SE(3)$ with rotations $R \in SO(3)$ and translations $b \in \R^3$, and let the component points of $P$ be $p^i \in \R^3$, then the rigid motion can be expressed as: 
    % \begin{equation}
    %     H \cdot P := \pmat{R p^1 + b \\ \vdots \\ R p^{N_i} + b} = \pmat{R p^1 \\ \vdots \\ R p^{N_i}} + \pmat{b \\ \vdots \\ b}\,.
    % \end{equation}  
    % We first focus on the rotating term. Any rotation of $P$ can be expressed as a linear combination of 9 linearly independent components $\bar{P}_{11},\cdots,\bar{P}_{33} \in \R^{3N_i}$, defined in the following. Write as $e_n \in \R^3$ the $n$-th unit vector, then:
    % \begin{align}
    %     \pmat{R p^1 \\ \vdots \\ R p^{N_i}}
    %     &= \sum_{n,m} R_{nm} \begin{pmatrix} (e_n^\top p^0) e_m\\ \cdots\\ (e_n^\top p^{N_i}) e_m \end{pmatrix} \\ &=: \sum_{n,m} R_{nm} \bar{P}_{nm}\,.
    % \end{align}
    % The vectors $\bar{P}_{nm}$ are linearly independent if $\dim\textnormal{span}\{p^i_{j,0} | i \in \{0,\cdots,N_i\}\} = 3$. If further $\dim\textnormal{span}\{R(t_k) | k \in \{0,\cdots,N_k\}\} = 9$ (which follows from the condition on $H(t_k)$), then the rotating component of a point cloud trajectory $P_{j,k}$  varies sufficiently to excite all $9$ linearly independent components $\bar{P}_{jk}$.

    % When each of the $N_j + 1$ trajectories corresponds to a unique static point cloud, we arrive at $9(N_j+1)$ linearly independent components to express $S$. 
    
    % Counting also the translations, these are the same for each point cloud since the dynamics~\eqref{eq:rigid-group-dynamics} are purely time-dependent. The translations can be expressed by three linearly independent vectors $\bar{P}_{b,1},\cdots, \bar{P}_{b,3}$ for all trajectories, bringing the maximum number of orthogonal components to $9(N_j+1)+3$. 
\end{proof}

\end{document}
